% !TEX root = ../../../main.tex
% Appendix entry: detection.seed_matches -- "Seeded search" (Appendix A).
\entryfield{Algorithm}
\code{_run} in \code{Adapters/detection_seed_matches.py}. The query $q$ is the \emph{exemplar}
side input's samples, $m = |q|$; it is resolved before the run by
\code{resolve_side_inputs} in \code{Working/side_inputs.py}, from a library exemplar (sliced fresh
from its recording's \code{.npy} by source file, channel and sample range), the chain's root signal,
or an earlier step's output.

\emph{The native search} (\emph{scales} = \texttt{1}, \code{match_exemplar}). The block calls
\code{stumpy.match}$(q, x, \text{max\_matches}=k, \text{max\_distance}=D)$ with $D = \infty$ when
\emph{max\_distance} is 0. stumpy computes the distance profile
$D_i = \lVert \hat q - \hat T_i \rVert_2$, the z-normalised Euclidean distance from $q$ to every
window $T_i = x_{i..i+m-1}$ (FFT-based MASS), then ranks and cuts greedily:
\begin{algorithmic}[1]
\State matches $\gets [\,]$; $e \gets \lceil \textit{exclusion} \times m \rceil$
\Loop
  \State $i \gets \arg\min_j D_j$
  \If{$D_i > D + 10^{-8}$ \textbf{or} $D_i$ is not finite \textbf{or} $|\text{matches}| = k$} \textbf{break} \EndIf
  \State append $(D_i, i)$; set $D_j \gets \infty$ for $i - e \le j \le i + e$
\EndLoop
\end{algorithmic}
So matches come out closest first, no two starting within $e$ samples of each other. stumpy reads
its zone from the global \code{stumpy.config.STUMPY_EXCL_ZONE_DENOM} (default 4, i.e.\ $m/4$); the
block sets it to $1/\textit{exclusion}$ for the call, under a module lock, and restores it. Each match
becomes the span $[i, i+m)$, labelled \texttt{match$r$} by rank $r$ and scored by $D_i$.

\emph{The scale bank} (any other \emph{scales}, \code{match_exemplar_bank}). For each factor $s$
the query is linearly resampled to $L = \mathrm{round}(m s)$ samples (\code{resample_to_length} in
\code{Working/distances.py}) and searched as above, with up to $k$ matches per length. Because a
z-normalised distance grows as the square root of the length, every distance is put on the native
footing, $d \sqrt{m/L}$ (and the cutoff divided by the same factor before the search). Matches of
\emph{different} lengths that overlap are then reduced by \emph{overlap}: \texttt{lowest} keeps the
smaller distance, \texttt{first} the earlier start, \texttt{all} keeps both (\code{reduce_overlaps}).
Two matches of one length are left to the exclusion zone. The survivors are sorted by distance and
the best $k$ kept; each span is as long as the copy that found it, labelled
\texttt{match$r$@$s$x}.

\entryfield{Assumptions}
\begin{itemize}
\item \emph{Units.} Z-normalisation removes offset and scale, so units do not matter, and neither
  does a constant offset between the exemplar's recording and this one.
\item \emph{Sampling rate.} The query is used sample for sample. Nothing compares the exemplar's
  rate with the span's: it is recorded in \code{meta["exemplar_fs"]} but not used. An exemplar
  from a 1\,Hz recording searched on a 10\,Hz span looks for a shape ten times faster than the one
  it was cut from.
\item \emph{Preprocessing.} A library exemplar is sliced from the stored channel as it is, while
  the span may have been filtered or detrended by earlier steps of the chain; the two are then
  compared after different treatment.
\item \emph{Missing values.} A NaN in the exemplar makes \code{stumpy.match} refuse (``Q contains
  illegal values''). A window of the span touching a NaN gets an infinite distance and is never a
  match.
\item \emph{Length.} The exemplar must have at least 3 samples and be no longer than the span.
\end{itemize}

\entryfield{Outputs}
A \texttt{SpanSet} of window-long spans, closest first, scored by distance and labelled as above;
the executor writes each to \code{detections}. \code{meta}: \code{m}, \code{exemplar_fs},
\code{n_matches}, \code{distances}, \code{exclusion} and \code{exclusion_samples}; with a bank also
\code{scales}, \code{per_scale} (matches per factor), \code{overlap} and \code{footing}.
Indices count from the start of the span.

In Discovery, the seed search (\code{Working/discovery/seeded_search.py}) uses this block's own
\code{match_exemplar} and \code{match_exemplar_bank} for its candidates, so the page and a run agree.
It adds what the block lacks: the same search on \term{surrogate}{surrogate} copies of the channel
(200 draws by default, phase randomisation unless Settings~\textrightarrow~Nulls says otherwise),
and a recommended cut, \code{recommended_cut}: the largest observed distance still below the
$\lfloor \alpha \times \text{draws}\rfloor$-th smallest null distance, so that chance would give at
most $\alpha$ matches per draw below it.

\entryfield{Cost}
One FFT distance profile per scale factor, about $O(n \log n)$ time and $O(n)$ memory: seconds on a
whole channel. The block declares no estimate and no span limit. Discovery multiplies it by
channels and null draws, and routes a seed search by a measured rate per channel
(\code{_seed_per_channel_s} in \code{webui/server/discovery.py}). Above the local ceiling,
Settings~\textrightarrow~Compute \& HPC's limit or ten minutes (\code{SEED_CEILING_S} $= 600$) when
none is saved, the search is not run in the bridge: \code{post_seed_slurm} writes a spec and an
\texttt{sbatch} script (\code{Working/discovery/seed_job.py}, \code{build_spec} and
\code{write_job}) that run the same candidates and null draws on the cluster, checkpointing and
resubmitting itself against the wall-clock limit, and the small result file is brought back through
\code{/api/discovery/seed/import}. \unverified{The page's \ui{Run seed search} button then reads
\ui{Create SLURM script}} (the endpoint's docstring).

\entryfield{Failure modes}
\begin{itemize}
\item \emph{Matches in noise.} With \emph{max\_distance} at 0 the block returns the $k$ best
  windows whatever they are. A ranked list is not evidence of recurrence; use Discovery's null and
  cut, or a calibrated \emph{max\_distance}.
\item \emph{The exemplar finds itself.} If the exemplar comes from the searched span (the root
  signal or a library entry on the same channel), \texttt{match0} is the exemplar at distance 0 and
  takes one of the $k$ places.
\item \emph{Overlapping matches.} At the default zone ($m/2$) two matches may share up to about
  half their samples, so one event can still be found twice, shifted. Set \emph{exclusion} to 1 or
  more to forbid any overlap. This is the overlapping-matches artefact of the older report tool
  (\code{Pipelines/motif_report/}); the block's default halves it rather than removing it. With a
  bank and \texttt{all}, matches of different lengths may overlap freely.
\item \emph{A cutoff that depends on length.} \emph{max\_distance} is an unnormalised
  z-normalised distance: it ranges up to $2\sqrt{m}$, so one value is strict for a long exemplar and
  lenient for a short one.
\item \emph{Sampling-rate mismatch} (above), unchecked.
\item \emph{Shared stumpy setting.} The zone is set by changing stumpy's global configuration. The
  lock serialises seed searches with one another, but a matrix profile or
  \blockref{detection.mp_motifs} run on another thread at that moment reads the changed zone.
\end{itemize}
The window-versus-event offset of the older tool does not arise in the same way: each span has
exactly the exemplar's footprint, so an exemplar cut to frame its event yields matches that frame
theirs the same way.

\entryfield{Validation}
\code{tests/test_adapter_seed_matches.py} plants a burst three times in noise and checks that the
exemplar finds itself at distance 0 and a second copy, closest first.
\code{tests/test_seed_exclusion_zone.py} checks that the zone defaults to $m/2$ and reaches stumpy.
\code{tests/test_library_seed_edges.py} checks that a bank finds a stretched shape at the right
length, the $\sqrt{m/L}$ footing and the overlap policies. No test measures how often matches on a
real recording are real events; that is the job of Discovery's null.

\entryfield{References}
\code{stumpy.match} \citep{law2019stumpy}; the distance is the matrix profile's
\citep{yeh2016matrix}. The $m/2$ zone is the project's choice (the spec's \S7.6).

\entryfield{Tests}
\code{tests/test_adapter_seed_matches.py}, \code{tests/test_seed_exclusion_zone.py},
\code{tests/test_library_seed_edges.py}; Discovery around it: \code{tests/test_discovery_seeded_search.py},
\code{tests/test_discovery_seed_job.py}, \code{tests/test_webui_discovery.py}.
